% =========================================================================
% abstract.tex
% =========================================================================

\begin{abstract}
Recent Key-Value (KV) cache quantization
methods, including TurboQuant, Quantized Johnson-Lindenstrauss (QJL), PolarQuant,
and Memory Caching, have demonstrated strong empirical results on large
language models with attention head dimension $\dhead \geq 128$.
Whether these methods generalize to smaller head dimensions remains
unexplored.
%
I present a systematic empirical evaluation at $\dhead = 64$ using
GPT-2 Medium (345M parameters) on a single NVIDIA RTX~4060 Laptop GPU.
I identify two failure modes at this scale: the QJL residual correction
introduces variance that degrades retrieval quality, and FibQuant angular
codebooks collapse to only four directions at four-bit precision.
%
I propose two corrections.
First, a simple fallback heuristic that disables QJL residual estimation below
$\dhead = 128$ achieves 67.4\% nearest-neighbor Recall@1 versus 38.8\%
for flat TurboQuant at $3.2{\times}$ compression.
Second, MC-TurboQuant combines Walsh-Hadamard Transform (WHT) vector
quantization with Gated Residual Memory (GRM), maintaining 97.5\%
average Needle-in-a-Haystack (NIAH) recall under quantization.
%
Empirical distortion remains within $1.83$--$2.46{\times}$ of the
Shannon lower bound across two to five bits.
A from-scratch training experiment reproduces the Memory Caching paper's
core finding at 15M-parameter scale (21.9 versus 81.9 perplexity).
%
This study is limited to a single model at $\dhead = 64$ without fused
CUDA kernels; cross-dimension generalization to $\dhead = 128$ was not
evaluated.
\end{abstract}
